\chapter{Ingestion des données dans HDFS}

\section{Arborescence HDFS}

Les données sont organisées en trois zones :

\begin{itemize}[leftmargin=1.5em]
  \item \texttt{/projet\_ia/raw/} : fichier CSV brut déposé depuis la
    source.
  \item \texttt{/projet\_ia/processed/} : sous-répertoires
    \texttt{offres/}, \texttt{metiers/} et \texttt{competences/},
    contenant les tables transformées au format Parquet.
  \item \texttt{/projet\_ia/output/} : exports et résultats finaux.
\end{itemize}

\section{Procédure d'ingestion}

L'ingestion est scriptée et exécutée automatiquement. Le script crée
l'arborescence HDFS, copie le CSV source dans \texttt{/projet\_ia/raw},
puis effectue les vérifications classiques (\texttt{hdfs dfs -ls},
\texttt{-du}, comparaison du nombre de lignes).

\section{Vérification}

Après dépôt dans HDFS :

\begin{itemize}[leftmargin=1.5em]
  \item \texttt{hdfs dfs -ls} : \texttt{/projet\_ia/raw/ai\_jobs\_market\_2025\_2026.csv},
    taille \texttt{383\,900} octets.
  \item \texttt{hdfs dfs -du -h} : \texttt{374.9 K}.
  \item Comparaison du nombre de lignes local vs HDFS
    (\texttt{wc -l} local vs \texttt{hdfs dfs -cat | wc -l}) :
    \textbf{1501 lignes des deux côtés} (1500 lignes de données + 1
    ligne d'en-tête). Le dépôt est intègre, aucune ligne n'a été perdue
    ou dupliquée pendant le transfert.
\end{itemize}
